\documentclass[11pt]{article}
\usepackage[margin=1in]{geometry}
\usepackage{amsmath,amssymb,amsthm,booktabs}
\usepackage[hidelinks]{hyperref}
\newtheorem{theorem}{Theorem}
\newtheorem{lemma}[theorem]{Lemma}
\newtheorem{proposition}[theorem]{Proposition}
\newtheorem{corollary}[theorem]{Corollary}
\newtheorem{conjecture}{Conjecture}
\renewcommand{\theconjecture}{13}
\theoremstyle{definition}
\newtheorem{definition}[theorem]{Definition}
\theoremstyle{remark}
\newtheorem{remark}[theorem]{Remark}
\newcommand{\R}{\mathbb{R}}
\newcommand{\C}{\mathbb{C}}
\newcommand{\Z}{\mathbb{Z}}
\newcommand{\floor}[1]{\lfloor #1\rfloor}
\usepackage{url}
\newcommand{\code}[1]{\path{#1}}

\title{Stability of the regular $n$-gon of point vortices at its bifurcation point:\\
completing the proof of the Montaldi--Tokieda conjecture}
\author{Gordon Speagle\\
Independent Researcher, Boise, ID, USA\\
\texttt{gspeagle@gmail.com}}
\date{Draft, 11 October 2026}

\begin{document}
\maketitle

\begin{abstract}
Montaldi and Tokieda studied a one-parameter family of point-vortex systems on the sphere, the plane and
the hyperbolic plane, and showed that a regular ring of $n\ge4$ identical vortices loses stability at an
explicit bifurcation value $b_n$ of the curvature parameter. They conjectured that the ring is still
stable \emph{at} $b_n$ (their Conjecture~13), in the sense of Lyapunov stability modulo rotations. They
proved it for all even $n$, where the sign of the decisive quartic coefficient is settled by a Maple
positivity check, and for $n=3,5,7,9,11$, where the odd-$n$ coefficients are evaluated numerically. We give
the first proof for odd $n\ge13$. Our argument is uniform in $n$, so it also re-proves the cases $4\le n\le11$
by a different route. Stability is Lyapunov stability modulo rotations about the ring's centre, obtained by
an energy--momentum argument on the momentum level set, and it holds for all small perturbations,
including those that change the momentum. The proof has three parts.
(i) A spectral computation. In log-polar coordinates the reduced Hessian at $b_n$ is block-diagonal and
independent of the geometry, and its kernel is exactly the critical Fourier mode.
(ii) Exact closed forms, for all $n$, for the quartic coefficient $c_4$ of the reduced Hamiltonian along
that kernel, after the transverse directions are eliminated on the momentum level set:
$c_4=(n-2)(m^4+40m^2+48m+16)/(192n(n-1))$ with $m=n-2$ for even $n$, and
$c_4=(n-1)(n-3)(n^2-9n+38)/(128n)$ for odd $n\ge5$. These replace Montaldi and Tokieda's Maple check (even
$n$) and numerical checks (odd $n\le11$). Positivity is immediate once the formulas are known; the work is
in deriving them.
(iii) A self-contained energy--momentum argument for degenerate (quartic) minima. It works for the
noncompact groups $SE(2)$ and $PSL(2,\R)$, and for the plane's affine momentum map.
The algebra that turns the computed series coefficients into the closed forms is checked in Lean~4.
We also compare our computation with some intermediate steps of the published argument; none of
this affects Montaldi and Tokieda's conclusions, which our results confirm. All verification so far has been by AI readers and
computer checks; no human specialist has yet reviewed the argument.
\end{abstract}

\tableofcontents

%======================================================================
\section{Introduction}\label{sec:intro}
%======================================================================

\paragraph{The family.} Montaldi and Tokieda~\cite{MT} put point vortices on the sphere, the plane and the
hyperbolic plane into one smooth family. In a uniform chart $z\in\C$ they use the surface $M_\lambda$ of
curvature $4\lambda$, with
\[
\sigma(z)=1+\lambda|z|^2,\qquad \Omega_\lambda=\frac{2}{\sigma^2}\,dx\wedge dy
\quad\text{(\cite[(1.8)]{MT})},\qquad
J_\lambda(z)=\frac1\sigma\big(z,|z|^2\big)\quad\text{(\cite[(1.9)]{MT})},
\]
and the pair interaction~\cite[(1.12)]{MT}
\[
G(z;w)=\log\frac{|z-w|^2}{(1+\lambda|z|^2)(1+\lambda|w|^2)} .
\]
For $\lambda>0$ this is the standard sphere (stereographic chart). For $\lambda=0$ it is the plane. For
$\lambda<0$ it is the hyperbolic plane with the ``sinh'' interaction, i.e.\ $G=2\log\sinh(d/2)$ up to
constants, where $d$ is hyperbolic distance. The symmetry group $G_\lambda$ is $SO(3)$, $SE(2)$ or
$PSL(2,\R)$ respectively. For $n$ identical unit vortices at distinct points $z_1,\dots,z_n$ we use
\begin{equation}\label{eq:H}
H=-\sum_{i<j}\log\frac{|z_i-z_j|^2}{\sigma_i\sigma_j},\qquad
\Omega=\sum_k\frac{2}{\sigma_k^2}dx_k\wedge dy_k,\qquad
J=(J_1,J_2,J_3)=\sum_k\Big(\frac{x_k}{\sigma_k},\frac{y_k}{\sigma_k},\frac{|z_k|^2}{\sigma_k}\Big),
\end{equation}
with $\sigma_k=\sigma(z_k)$. This is the Hamiltonian of~\cite{MT} without the positive factor $\kappa^2/4\pi$.
That factor rescales time and does not affect any statement below.

\paragraph{Rings and the bifurcation point.} The regular ring $z_k=r_0e^{2\pi ik/n}$ is a relative
equilibrium for every $r_0$ and $\lambda$. By scaling we may take $r_0=1$, so the only parameter is
$b=\lambda r_0^2\in(-1,1]$. Montaldi--Tokieda~\cite[Thm.~6]{MT} prove that for $n\ge4$ the ring is
stable for $\lambda r_0^2<b_n$, where $b_n$ is the unique root in $(-1,1]$ of
\begin{equation}\label{eq:bn}
\frac{1+b_n^2}{(1+b_n)^2}=\frac{1}{2(n-1)}\Big\lfloor\frac{n^2}{4}\Big\rfloor
\qquad\text{(\cite[(3.6)]{MT})}.
\end{equation}
Above $b_n$ the ring is spectrally unstable, and a pitchfork bifurcation occurs. Explicitly,
\[
b_n=\frac{4(n-2)\sqrt{n-1}-n^2}{n^2-8n+8}\ \ (n\text{ even}),\qquad
b_n=\frac{4\sqrt{n-3}-(n+1)}{n-7}\ \ (n\text{ odd},\ n\ne7),\qquad b_7=0 .
\]
So $b_4=2-\sqrt3$, $b_5$ and $b_6=9-4\sqrt5$ are positive (sphere). The value $b_7=0$ is the plane, where
this is the Thomson heptagon. For $n\ge8$, $b_n<0$ (hyperbolic plane); for example $b_8=-(8-3\sqrt7)$.
Throughout, $|b_n|<1$.

\paragraph{The conjecture.} Page numbers for~\cite{MT} refer to arXiv:1210.5662v1. Montaldi and Tokieda
define stability as follows~\cite[p.~10]{MT}:
``The precise sense of stability we are adopting is like the ordinary one of Lyapunov stability, but in
terms of $G_\mu$-invariant open sets, where $\mu$ is the momentum value at the trajectory and $G_\mu$ is
the stabilizer of $\mu$: in detail, for every $G_\mu$-invariant neighbourhood $V$ of the trajectory, there
exists a $G_\mu$-invariant neighbourhood $U\subseteq V$ such that every trajectory starting in $U$ remains
in $V$ for all time.'' Before Theorem~6 they add~\cite[p.~15]{MT}: ``Stability means Lyapunov stability
modulo rotations (same as orbital stability in our situation).'' Then, on p.~23:

\begin{quote}
\textbf{Conjecture 13.} \emph{For all values of $n$, the $n$-gon is stable at the bifurcation point.}\\
Above, we have proved this for all even $n$ and for $n = 3, 5, 7, 9, 11$.
\end{quote}

\paragraph{The notion of stability.} Throughout, ``stable'' means Lyapunov stability modulo rotations
about the ring's centre, in the sense of~\cite[p.~10, p.~15]{MT} just quoted (Theorem~\ref{thm:main} states it
precisely). We obtain it by the energy--momentum method: the reduced Hamiltonian on the momentum level set
$J^{-1}(\mu)/K$ has a strict, but degenerate (quartic), local minimum, and Section~\ref{sec:nonlinear} turns
this into Lyapunov stability modulo $K$ in the full phase space.

\paragraph{What was known.} Montaldi and Tokieda~\cite{MT} proved Conjecture~13 for all even $n$ and for
$n=3,5,7,9,11$, as quoted above. For even $n$~\cite[\S3.5.1, pp.~20--21]{MT}, the cubic term vanishes on the
slice (their Lemma~10), and the $t^4$ coefficient along the critical mode is a positive multiple of a
bracket $T$. Substituting the two roots of the threshold relation gives two values $T_1,T_2$, and closed
forms for $T_1+T_2$ and $T_1T_2$, both manifestly positive for $m=n-2\ge0$, are obtained ``using Maple!''
(p.~21). For odd $n$~\cite[\S3.5.2, pp.~22--23]{MT}, the criterion involves three
coefficients $\beta,\gamma,\delta$. The coefficient $\beta$ has a closed form, while ``the awkward
trigonometric expressions prevented us (and Maple) from reaching closed forms for $\gamma$ and $\delta$''
(p.~23). The coefficients are evaluated numerically for $n=5,9,11$. For $n=7$ they are found with Maple
as rational multiples of $1/\pi r_0^k$, which are ``correct to a high degree of precision---but we have no
proof that they are rational multiples of $1/\pi r_0^k$''; the sign of $\beta\delta-\gamma^2$ is then
certain. Their case
$n=3$ rests on their statement (p.~15) that ``the ring of $n=3$ vortices is always nonlinearly stable''. So
odd $n\ge13$ (all in the hyperbolic plane) remained open. The case $n=7$ is the Thomson heptagon. Its
nonlinear stability was proved by Kurakin and Yudovich~\cite{KY} (we have read only the abstract) and by
Schmidt~\cite{Schmidt}; Montaldi--Tokieda recover it~\cite[Cor.~7]{MT}. For the sphere at the critical
latitude ($n=4,5,6$), Borisov and Kilin~\cite{BK} wrote that the problem ``requires separate analysis with
respect to nonlinear terms and remains unsolved to the present day''. Kurakin~\cite{Kurakin03,Kurakin04}
later gave necessary and sufficient conditions for nonlinear stability of regular vortex systems on the
sphere. We have not read those papers and make no claim about their content.

\paragraph{Contributions.}
\begin{enumerate}
\item The first proof of Conjecture~13 for odd $n\ge13$. Together with~\cite{MT} this settles the
conjecture for all $n$. Our proof (Theorem~\ref{thm:main}) is uniform for all $n\ge4$, so it gives an
independent proof of the cases $4\le n\le11$ already proved in~\cite{MT}. Stability is in
Montaldi--Tokieda's sense, and it holds for perturbations of arbitrary (nearby) momentum.
\item Exact closed forms for the reduced quartic coefficient $c_4$, for every even and every odd $n$
(Theorems~\ref{thm:even} and~\ref{thm:odd}). They replace the Maple positivity check of~\cite{MT} for even
$n$ and their numerical evaluations for odd $n=5,7,9,11$. Once the formulas are known, positivity is
immediate (Remark~\ref{rem:positivity}); the content is the derivation.
\item A self-contained nonlinear-stability argument for degenerate minima that covers the noncompact groups
involved, including the plane with its affine momentum map (Section~\ref{sec:nonlinear}).
\item A Lean~4 check of the algebra that leads from the computed series coefficients to the closed forms
and to Table~\ref{tab:c4} (Section~\ref{sec:limits}, which states exactly what is and is not formalised).
\item A comparison with intermediate steps of~\cite{MT}, given with page references so that the reader can
check our reading (Section~\ref{sec:remarks}). None of it affects any conclusion of~\cite{MT}.
\end{enumerate}

%======================================================================
\section{Main result}\label{sec:main}
%======================================================================

\begin{theorem}\label{thm:main}
Let $n\ge4$. Consider $n$ identical point vortices on $M_\lambda$ with Hamiltonian~\eqref{eq:H}, in a
regular ring at the bifurcation point $\lambda r_0^2=b_n$. This means the sphere for $n=4,5,6$, the plane
for $n=7$ and the hyperbolic plane for $n\ge8$. Let $K\cong SO(2)$ be the group of rotations about the ring
centre, so that $K=G_\mu$ for $\mu=J(\text{ring})$. Then the ring is Lyapunov stable modulo $K$. That is,
for every $\varepsilon>0$ there is $\delta>0$ such that every solution with
$\operatorname{dist}(x(0),K\cdot x_0)<\delta$ exists for all $t\in\R$ and satisfies
$\operatorname{dist}(x(t),K\cdot x_0)<\varepsilon$. No restriction is placed on the momentum of $x(0)$.
\end{theorem}

Since $K$ is compact, $K$-invariant neighbourhoods of the (compact) relative-equilibrium orbit $K\cdot x_0$
are cofinal among all of its neighbourhoods. So Theorem~\ref{thm:main} is exactly stability in the sense
of~\cite[p.~10]{MT} quoted above. It also implies stability modulo the full group $G_\lambda$.

\begin{corollary}
Conjecture~13 of~\cite{MT} holds for every $n\ge4$. Together with the case $n=3$, which~\cite{MT} state is
always nonlinearly stable~\cite[p.~15]{MT}, it holds for all $n$.
\end{corollary}

For $n=3$, equation~\eqref{eq:bn} has its only root at $b=1$ (the equator). There the momentum value is
fixed by all of $SO(3)$, the hypothesis $|b|<1$ used throughout fails, and the critical mode is $m=1$,
which meets the momentum directions. Consistently with this, the odd-$n$ formula of
Theorem~\ref{thm:odd} gives $c_4=0$ at $n=3$, so it cannot decide that case, and its derivation assumes
$n\ge5$ in any case. We rely on Montaldi--Tokieda's statement for $n=3$ and do not re-prove it. For
$n\ge4$ the new content is odd $n\ge13$; for $4\le n\le11$, Theorem~\ref{thm:main} is an independent proof
of results already in~\cite{MT}.

\paragraph{Structure of the proof.} Section~\ref{sec:setup} shows that the reduced Hessian at $b_n$ is
positive semidefinite, with kernel equal to the radial part of the critical Fourier mode
(Proposition~\ref{prop:spectral}). Section~\ref{sec:quartic} computes the quartic coefficient of the
reduced Hamiltonian along this kernel, after the transverse directions are eliminated, and shows it is
positive (Theorems~\ref{thm:even} and~\ref{thm:odd}). Section~\ref{sec:nonlinear} turns ``degenerate strict
minimum of the reduced Hamiltonian'' into Lyapunov stability modulo $K$ in the full phase space
(Theorem~\ref{thm:abstract}).

%======================================================================
\section{Setup: Fourier blocks and the reduced space}\label{sec:setup}
%======================================================================

\subsection{Log-polar coordinates and the Fourier blocks}

Write $z_k=e^{s_k+i(\phi_k+\theta_k)}$ with $\phi_k=2\pi k/n$, so the ring is $s=\theta=0$. Then
\[
\log|z_j-z_k|^2=s_j+s_k+\log\big(2\cosh(s_j-s_k)-2\cos(\phi_j-\phi_k+\theta_j-\theta_k)\big),
\]
which does not involve $\lambda$. All dependence on $\lambda$ is in one-body terms that depend on the
$s_k$ alone. These are $(n-1)\sum_k\log(1+b\,e^{2s_k})$ in $H$, and
$J_3=\sum_k g(e^{2s_k})$ with $g(u)=u/(1+bu)$.

Let $F=H-\omega_0J_3$, with $\omega_0$ chosen so that $dF(\text{ring})=0$. With
$D_\phi=2-2\cos\phi$, the pair logarithm has no linear $\Delta s$ term and no
$\Delta s\,\Delta\theta$ term at second order; its quadratic part is $(\Delta s^2-\Delta\theta^2)/D_\phi$.
The Fej\'er-type identity
\begin{equation}\label{eq:fejer}
\sum_{l=1}^{n-1}\frac{1-\cos(2\pi ml/n)}{1-\cos(2\pi l/n)}=m(n-m)\qquad(0\le m\le n)
\end{equation}
(equivalent to~\cite[(2.4)]{MT}) diagonalises the pair part on each Fourier mode. Its proof is one line:
with $w=e^{2\pi il/n}$ the summand is $|1+w+\dots+w^{m-1}|^2$, then sum over $l$.

The one-body part of $F$ is $\sum_k\varphi(s_k)$ with
$\varphi(s)=-(n-1)s+(n-1)\log(1+be^{2s})-\omega_0\,e^{2s}/(1+be^{2s})$. The condition $\varphi'(0)=0$ gives
$\omega_0=-(n-1)(1-b^2)/2$, and then
\[
C_1:=\varphi''(0)=\frac{4(n-1)b}{(1+b)^2}+\frac{(n-1)(1-b^2)}2\cdot\frac{4(1-b)}{(1+b)^3}
=\frac{2(n-1)(1+b^2)}{(1+b)^2}.
\]
Let $V_m$ ($0\le m\le\floor{n/2}$) be the real Fourier spaces, spanned by
$(\cos m\phi_k,\sin m\phi_k)$ in $s$ and in $\theta$. The space $V_m$ has dimension $4$ for
$0<m<n/2$, and dimension $2$ for $m=0$ and for $m=n/2$.

\begin{proposition}[Fourier blocks]\label{prop:blocks}
For every $b\in(-1,1)$, $d^2F(\text{ring})$ is block-diagonal over the $V_m$, with no radial--tangential
coupling. Its eigenvalues are
\[
a_m=C_1-m(n-m)\ \ \text{(radial)},\qquad b_m=m(n-m)\ \ \text{(tangential)},
\]
with respect to the Euclidean structure $\sum_k(ds_k^2+d\theta_k^2)$. At $b=b_n$ we have
$C_1=\floor{n^2/4}=\max_m m(n-m)$, which is~\eqref{eq:bn}.
\end{proposition}

\noindent
This is~\cite[(2.6)]{MT} up to normalisation. We have $a_m\ge0$ at $b_n$, with equality iff
$m=\floor{n/2}$, because $m(n-m)$ is strictly increasing for $1\le m\le n/2$. For odd $n$, the modes
$(n\pm1)/2$ form the single real space $V_{(n-1)/2}$.

\subsection{The reduced space}

Let $x_0$ be the ring and $\mu=J(x_0)=(0,0,\mu_3)$ with $\mu_3=n/(1+b)$. The stabiliser $G_\mu$ is the group
$K$ of rotations about the origin, because $|b|<1$ excludes the equator. Let $t$ denote the rotation
direction, $\theta_k=1$ for all $k$.

\begin{enumerate}
\item[(a)] $dJ_3=2(1+b)^{-2}\sum_kds_k$ is nonzero exactly on the radial part of $V_0$.
\item[(b)] Let $J_+=J_1+iJ_2$, and let $R$ be the rotation by $2\pi/n$ composed with the relabelling
$k\mapsto k+1$. Then $R$ fixes the ring and $J_+\circ R=e^{2\pi i/n}J_+$. So $dJ_+(x_0)$ vanishes on $V_m$
for $m\ne1$. Explicitly $dJ_+=\sum_ke^{i\phi_k}(A\,ds_k+iB\,d\theta_k)$, with $A=(1-b)/(1+b)^2$ and
$B=1/(1+b)$. It has real rank $2$ on $V_1$; let $V_1'=V_1\cap\ker dJ_+$, which has dimension $2$.
\item[(c)] The other two generators $\xi_1,\xi_2$ of $G_\lambda$ are rotations about horizontal axes,
translations, or boosts. They span a $2$-plane in $V_1$ that is \emph{not} in $\ker dJ$, because
$\det\big(dJ_i(\xi_j)\big)_{i,j=1,2}=\{J_1,J_2\}(x_0)^2$ and
\begin{equation}\label{eq:bracket}
\{J_1,J_2\}=\tfrac n2-\lambda J_3,\qquad \{J_1,J_2\}(x_0)=\frac{n(1-b)}{2(1+b)}\ne0 .
\end{equation}
For $\lambda=0$ the bracket is the constant $n/2$. This is the central term of the affine (non-equivariant)
momentum map of $SE(2)$. Our argument never uses equivariance, only that $J_1,J_2,J_3$ are Hamiltonians
for the generators and the value~\eqref{eq:bracket}. For $\lambda\ne0$, $J-(0,0,n/(2\lambda))$ is equivariant.
\end{enumerate}

Hence $\ker dJ(x_0)=\R t\oplus V_1'\oplus V_2\oplus\dots\oplus V_{\floor{n/2}}$. The $G_\lambda$-orbit meets
$\ker dJ$ only in $\R t$. The reduced space $M_\mu=J^{-1}(\mu)/K$ is a smooth manifold near $[x_0]$ ($K$
acts freely on labelled rings), and
\[
\dim M_\mu=2n-4 .
\]
Three dimensions are removed by constraints (the radial part of $V_0$, and two directions in $V_1$), and one by
the gauge direction $t$.

\subsection{Gap (i): spectral positivity off the kernel}

Let $h_\mu$ be the reduced Hamiltonian. On $J^{-1}(\mu)$ we have $H=F+\omega_0\mu_3$. For any curve $c$
in $J^{-1}(\mu)$ through $x_0$, $(F\circ c)''(0)=d^2F(c',c')+dF(c'')=d^2F(c',c')$, because $dF(x_0)=0$. So
$d^2h_\mu([x_0])$ is $d^2F(x_0)$ restricted to $\ker dJ(x_0)$, with $t$ in its radical.

\begin{proposition}\label{prop:spectral}
Let $n\ge4$ and $\lambda r_0^2=b_n$, in any of the three geometries. The Hessian of $h_\mu$ at $[x_0]$ is
positive semidefinite. Its kernel is exactly the radial part of $V_{\floor{n/2}}$, of dimension $1$ for
even $n$ and $2$ for odd $n$. It is positive definite on the complement
$E=V_1'\oplus\bigoplus_{2\le m\le\floor{n/2}}(\text{tangential }V_m)\oplus\bigoplus_{2\le m<\floor{n/2}}
(\text{radial }V_m)$.
\end{proposition}

\begin{proof}
By Proposition~\ref{prop:blocks} the restriction is block-diagonal.
\emph{Mode $1$:} on all of $V_1$, $d^2F=\operatorname{diag}(a_1I,b_1I)$ with $b_1=n-1>0$ and
$a_1=\floor{n^2/4}-(n-1)>0$ for $n\ge4$. So $d^2F$ is positive definite on every subspace of $V_1$,
in particular on $V_1'$. Nothing about $A$ or $B$ is needed.
\emph{Modes $m\ge2$:} $b_m>0$, and $a_m>0$ unless $m=\floor{n/2}$, where $a_m=0$.
\emph{Mode $0$:} its radial part is removed by $J_3=\mu_3$ and its tangential part is the gauge direction.
\end{proof}

\noindent
The smallest nonzero eigenvalue (in the $(s,\theta)$-Euclidean normalisation) is $1$ for even $n\ge6$
($a_{n/2-1}=1$) and $2$ for odd $n\ge7$ ($a_{(n-3)/2}=2$). For $n=4,5$ the next mode is $m=1$, where
$V_1'$ mixes radial and tangential parts; the values are $3/2$ and $8/3$. An independent 40-digit computation of the full
$2n\times2n$ Hessian, compressed onto $\ker dJ\cap t^\perp$, agrees with these predictions for $n=4,\dots,15$
to $5\cdot10^{-39}$. A bordered-inertia check at $n=5,7,\dots,25,31$ finds exactly the predicted kernel and
no negative eigenvalue (Section~\ref{sec:limits}).

%======================================================================
\section{The effective quartic on the level set}\label{sec:quartic}
%======================================================================

\subsection{Why the level set, and not a linear slice}

Let $K_0$ be the kernel in Proposition~\ref{prop:spectral}, with coordinates $(a,b)\in K_0\times E$ on
$M_\mu$. Let $\beta(a)$ be the slaving map, the critical point of $b\mapsto h_\mu(a,b)$, and let
$g(a)=h_\mu(a,\beta(a))-h_\mu(0)$. Because the Hessian vanishes on $K_0$, the leading term of $g$ is of
order $4$ (the cubic term vanishes by symmetry; see below). It is built from three ingredients: the quartic
derivative of $H$ along $a$, the cubic coupling of $a^2$ to the slaved $O(|a|^2)$ displacement, and the
quadratic energy of that displacement.

The slaved displacement must keep the curve on $J^{-1}(\mu)$. Moving a distance $\varepsilon$ along a
kernel vector changes $J_3$, and through $J_+$ also the mode-$1$ momenta, at order $\varepsilon^2$. So a
curve in the level set has an $O(\varepsilon^2)$ component along directions \emph{not} in
$\ker dJ(x_0)$, chiefly the radial part of $V_0$. This component has Hessian weight $a_0=C_1\ne0$ and
couples to $a^2$ at cubic order, so it contributes at order $\varepsilon^4$.

When the Hessian is nondegenerate, $O(\varepsilon^4)$ is negligible next to the quadratic term, and the
curvature of the level set is invisible. When the Hessian is degenerate, $\varepsilon^4$ is the leading
order, and the curvature of $J^{-1}(\mu)$ enters the decisive coefficient. A computation on the linear
slice $x_0+N$ ($N$ a symplectic normal space) therefore does not in general give the quartic coefficient
of the reduced Hamiltonian. Section~\ref{sec:remarks} quantifies this.

\begin{definition}\label{def:c4}
Let $e\in\R^{2n}\cong\C^n$ be the Euclidean-unit radial kernel vector in the chart:
$e_k=n^{-1/2}(-1)^k\hat z_k$ for even $n$, and $e_k=(2/n)^{1/2}\cos(l\phi_k)\hat z_k$ with
$l=(n-1)/2$ for odd $n$. Here $\hat z_k$ is the outward unit vector at the $k$-th vortex. For odd $n$, let
$e'$ be the same with $\sin$ in place of $\cos$. Define
\begin{multline*}
f_{\rm eff}(\varepsilon)=\min\big\{H(q)-H(x_0):\ J(q)=\mu,\ \langle q-x_0,e\rangle=\varepsilon,\\
\langle q-x_0,e'\rangle=0\ (n\text{ odd}),\ \langle q-x_0,t\rangle=0,\ q\text{ near }x_0\big\}.
\end{multline*}
Then $f_{\rm eff}(\varepsilon)=c_4\varepsilon^4+O(\varepsilon^5)$, and $c_4$ is the \emph{level-set
effective quartic}.
\end{definition}

\noindent
$f_{\rm eff}$ is $g$ restricted to a ray in $K_0$, written in a different chart. Positivity of the quartic
term does not depend on the chart: a change of chart acts on $K_0$ by an invertible linear map to leading
order. By symmetry, $g$ has no cubic term and its quartic term is $c_4|a|^4$:
\begin{itemize}
\item For even $n$, the generator of $\Z_n$ acts by $-1$ on $K_0=\text{radial }V_{n/2}$, so $g$ is even.
\item For odd $n$, $\Z_n$ acts on $K_0\cong\C$ by $a\mapsto e^{2\pi il/n}a$ with $\gcd(l,n)=1$. The
monomial $a^p\bar a^q$ is invariant iff $n\mid p-q$. For $p+q\le4<n$ ($n\ge5$) this forces $p=q$. So there
is no cubic, and the only quartic is $|a|^4$, as in~\cite[Prop.~12]{MT}.
\end{itemize}
Since $h_\mu$, $J$ and the Fourier splitting are $\Z_n$-equivariant, $g$ is $\Z_n$-invariant.

\subsection{Even $n$: symmetric criticality and product formulas}

\begin{theorem}\label{thm:even}
For even $n=2M\ge4$ at $b=b_n$,
\[
c_4=\frac{(n-2)\,(m^4+40m^2+48m+16)}{192\,n\,(n-1)},\qquad m=n-2 .
\]
In particular $c_4>0$, and $c_4=\tfrac14,\tfrac{23}{30},\tfrac{95}{56},\tfrac{49}{15},\tfrac{755}{132}$ for
$n=4,6,8,10,12$.
\end{theorem}

\begin{proof}
\emph{Symmetric criticality.} The kernel vector $e$ is fixed by the subgroup $D_M\subset D_n$ generated by
rotation by $2\pi/M$ and the reflection fixing vortex $0$. For fixed $\varepsilon$ the slaved problem is a
strictly convex minimisation near $x_0$ (Proposition~\ref{prop:spectral} and the implicit function
theorem). It is $D_M$-equivariant, so its unique minimiser is $D_M$-invariant. The level set and the
constraints are also $D_M$-invariant. Hence the minimiser lies in $\operatorname{Fix}(D_M)\cap J^{-1}(\mu)$.
Up to rotation, this set is the curve of configurations made of two staggered regular $M$-gons,
\[
\{\rho_+\eta^j\}_{j<M}\cup\{\rho_-\zeta\eta^j\}_{j<M},\qquad \eta=e^{2\pi i/M},\ \zeta=e^{i\pi/M},
\]
with radii constrained by $J_3=\mu_3$. ($J_1=J_2=0$ holds automatically.) So $f_{\rm eff}$ is, to all orders,
the restriction of $H$ to this curve.

\emph{Product formulas.} Put $u_\pm=\rho_\pm^2$. Within one $M$-gon,
$\sum_{i<j}\log|z_i-z_j|^2=\binom M2\log u_\pm+\text{const}$. Between the two,
\[
\sum_{j,k}\log|\rho_+\eta^j-\rho_-\zeta\eta^k|^2=M\log\Big|\prod_{\omega^M=-1}(\rho_+-\rho_-\omega)\Big|^2
=2M\log\big(u_+^{M/2}+u_-^{M/2}\big),
\]
using $\prod_{\omega^M=-1}(x-y\omega)=x^M+y^M$. With $\sigma_\pm=1+\lambda u_\pm$, this gives
\[
H=-\tbinom M2(\log u_++\log u_-)-2M\log\big(u_+^{M/2}+u_-^{M/2}\big)+(n-1)M(\log\sigma_++\log\sigma_-)
+\text{const}.
\]
\emph{Expansion on the constraint.} The constraint is $g(u_+)+g(u_-)=2g(1)$ with $g(u)=u/(1+\lambda u)$.
Parametrise it by $g(u_\pm)=g(1)\pm\tau$. The swap $\tau\mapsto-\tau$ is a rotation by $\pi/M$, so $H$ is
even in $\tau$: $H=H_0+h_2\tau^2+h_4\tau^4+O(\tau^6)$. Direct expansion gives
\[
h_2=-\frac M4(1+\lambda)^2\big[M^2(1+\lambda)^2-(4M-2)(1+\lambda^2)\big].
\]
This vanishes exactly when $(1+\lambda^2)/(1+\lambda)^2=M^2/(2(2M-1))=n^2/(8(n-1))$, which is~\eqref{eq:bn}
for even $n$. The expansion also gives
\[
h_4=\frac{M(1+\lambda)^4}{96}\Big[M^4(1+\lambda)^4-16M^2(1+\lambda)(1+\lambda^3)+24(2M-1)(1+\lambda^4)\Big].
\]
The $e$-component of the displacement is $\varepsilon=\tfrac{\sqrt n}{2}(\rho_+-\rho_-)+O(\rho-1)^2$, and
$\tau=\tfrac{2}{(1+\lambda)^2}\,\tfrac{\rho_+-\rho_-}2+O(\cdot)^2$. Since $h_2=0$, only the leading linear
relation matters, and
\[
c_4=h_4\Big(\frac{2}{(1+\lambda)^2}\Big)^4\frac1{n^2}.
\]
Reducing modulo the quadratic relation~\eqref{eq:bn} gives the stated rational function of $n$. The
reduction was done symbolically in $M$ and $\lambda$, by two independent scripts
(Section~\ref{sec:limits}). Positivity is clear: $n-2>0$ and $m^4+40m^2+48m+16>0$.
\end{proof}

\subsection{Odd $n$: canonical coordinates, root-of-unity sums and residues}

\begin{theorem}\label{thm:odd}
For odd $n\ge5$ at $b=b_n$,
\[
c_4=\frac{(n-1)(n-3)(n^2-9n+38)}{128\,n}.
\]
In particular $c_4>0$, because $n^2-9n+38$ has discriminant $-71$. Also $c_4(7)=9/14$.
\end{theorem}

\begin{proof}
\emph{Canonical coordinates.} Put $w=|z|^2/\sigma$ and $\vartheta=\arg z$. Then
$\Omega=\sum_kdw_k\wedge d\vartheta_k$, because $2\,dx\wedge dy/\sigma^2=du\wedge d\vartheta/\sigma^2$ with
$u=|z|^2$, and $dw=du/\sigma^2$. In these coordinates
\[
J_3=\sum_kw_k,\qquad J_1+iJ_2=\sum_k\sqrt{P(w_k)}\,e^{i\vartheta_k},\qquad P(w)=w(1-\lambda w),
\]
\[
\frac{|z_i-z_j|^2}{\sigma_i\sigma_j}=w_i+w_j-2\lambda w_iw_j-2\sqrt{P(w_i)P(w_j)}\cos(\vartheta_i-\vartheta_j).
\]
So $H$ has no one-body term, and $J_3$ is linear. The ring is $w_k=w_0=1/(1+\lambda)$,
$\vartheta_k=\phi_k$.

\emph{Ansatz.} Write $w_k=w_0+\delta_k$ and $\vartheta_k=\phi_k+\psi_k$. On $J^{-1}(\mu)$ we have
$\sum\delta_k=0$ \emph{exactly}, because $J_3$ is linear. Fix the gauge by $\sum\psi_k=0$, and take the
kernel direction $\delta_k=A\cos(l\phi_k)$, $l=(n-1)/2$. Since $2l\equiv-1\pmod n$, $a^2$ lies in Fourier
modes $0$ and $\pm1$, so only these modes are slaved at order $A^2$. Mode $0$ is removed, as just said. The
reflection $k\mapsto-k$ fixes the kernel vector, so the slaved mode-$1$ part has the form
\[
\delta_k=A\cos(l\phi_k)+A^2p\cos\phi_k,\qquad \psi_k=A^2s\sin\phi_k .
\]
Corrections of order $A^3$ do not affect the $A^4$ coefficient. The partial gradient $\partial_bh_\mu$ at
the $O(A^2)$-truncated slaved point is $O(A^3)$, and pairing it with an $O(A^3)$ correction gives
$O(A^6)$. Modes $m\ge2$ have no $a^2$ coupling, so their slaved amplitude is $O(A^3)$ and contributes
$O(A^6)$ as well.

\emph{Constraint.} $J_2=0$ holds by reflection symmetry. At order $A^2$, the condition $J_1=0$ reads
\[
\tfrac n2A^2\big[Q'(w_0)p-Q(w_0)s+\tfrac14Q''(w_0)\big]=0,\qquad Q=\sqrt P,
\]
with $Q(w_0)=1/(1+\lambda)$, $Q'(w_0)=(1-\lambda)/2$ and $Q''(w_0)=-(1+\lambda)^3/4$. Here we used
$\sum_k\cos^2(l\phi_k)\cos\phi_k=n/4$. The inhomogeneous term $\tfrac14Q''$ is the curvature of the level
set, which in these coordinates sits in $J_1$ rather than $J_3$.

\emph{Sums over roots of unity.} Expanding $H$ along the ansatz requires sums over $k$ of rational
functions of $e^{i\phi}$ and $e^{il\phi}$ for offsets $\phi=2\pi m/n$. Put $t=e^{i\pi m/n}$. Then
$e^{i\phi}=t^2$ and $e^{il\phi}=(-1)^m/t$, and $\sum_{m=1}^{n-1}$ becomes half the sum over the $2n$-th
roots of unity $t\ne\pm1$. For each sign $\varsigma=t^n=\pm1$, the sum $\sum_{t^n=\varsigma}R(t)$ is
evaluated by residues of $R(t)\,nt^{n-1}/(t^n-\varsigma)$. The finite poles of $R$ are only at $t=\pm1$
(and possibly $0$). Every piece has Laurent degree $\le0$ at $0$ and at $\infty$, so the residue at $0$
vanishes and the residue at $\infty$ is $-n$ times the constant term of $R$ at $\infty$. The residues at
$\pm1$ are computed from binomial expansions of $t^n$ about $\pm1$, which are polynomials in $n$. Every
sum is therefore an explicit rational function of $n$ and $\lambda$, valid for all odd $n\ge5$. This is the
step that replaces the numerical evaluation in~\cite{MT}.

\emph{Result.} The $A^2$ coefficient of $H$ along the kernel is
\[
f_2=-\frac{n(n-1)(1+\lambda)^2}{64}\big[(n-7)(1+\lambda^2)+2(n+1)\lambda\big],
\]
which vanishes exactly at $b_n$; for odd $n$,~\eqref{eq:bn} reads $(1+\lambda^2)/(1+\lambda)^2=(n+1)/8$.
The $A^3$ coefficient vanishes. The $A^4$ coefficient is
$E(p,s)=E_{00}+E_{10}p+E_{01}s+E_{20}p^2+E_{02}s^2$ (there is no $ps$ term), with
\begin{align*}
E_{20}&=\tfrac1{16}n(n-1)(1-\lambda)^2(1+\lambda)^2, &
E_{02}&=\tfrac14n(n-1), \\
E_{10}&=\tfrac1{64}n(n-1)(1-\lambda)(1+\lambda)^3\big[(n-5)(1+\lambda^2)+2(n-1)\lambda\big],&
E_{01}&=\tfrac1{64}n(n-1)^2(1+\lambda)^4,
\end{align*}
\begin{multline*}
E_{00}=\tfrac{n(n-1)(1+\lambda)^4}{8192}\Big[(n^3+n^2-69n+291)(1+\lambda^4)\\
+4(n^3+n^2-21n-13)(\lambda+\lambda^3)+(6n^3+6n^2-30n+82)\lambda^2\Big].
\end{multline*}
Eliminate $s$ with the $J_1$ constraint. The coefficient of $p^2$ in what remains is
$a_2=\tfrac18n(n-1)(1-\lambda)^2(1+\lambda)^2>0$, so the slaved $p$ is a minimum. Normalising
$\delta_k=A\cos(l\phi_k)$ to the unit vector $e$ gives $A=\varepsilon\,(2/n)^{1/2}\,2/(1+\lambda)^2$, so
\[
c_4=\frac{64}{n^2(1+\lambda)^8}\min_p E\big(p,s(p)\big).
\]
Reducing modulo the threshold relation $(n-7)(1+\lambda^2)+2(n+1)\lambda=0$ gives the stated formula for
$n\ne7$. The denominator does not vanish modulo the relation unless $n=7$. Direct substitution
$\lambda=0$, $n=7$ gives $9/14$, which agrees with the formula.
\end{proof}

\begin{remark}[Where the difficulty lies]\label{rem:positivity}
Once the closed forms are known, $c_4>0$ is immediate. For even $n$ every term of
$(n-2)(m^4+40m^2+48m+16)$ is positive for $m=n-2\ge2$. For odd $n$, $n^2-9n+38$ has discriminant $-71$, and
$(n-1)(n-3)>0$ for $n\ge5$. The content of Theorems~\ref{thm:even} and~\ref{thm:odd} is the derivation of
$c_4$, in particular the evaluation of the root-of-unity sums for general $n$, which~\cite[p.~23]{MT} could
not carry out in closed form. The odd formula vanishes at $n=3$, because of the factor $n-3$. This is
consistent with $n=3$ being a different situation (bifurcation at the equator, $b=1$), which needs, and
in~\cite{MT} receives, separate treatment; see the remark after the Corollary in Section~\ref{sec:main}.
\end{remark}

\begin{remark}[What $c_4$ is, for odd $n$]\label{rem:reduced}
For odd $n$, $c_4$ is not the quartic coefficient of $H$ along the kernel (the analogue of Montaldi and
Tokieda's $\delta$). It is the quartic coefficient of the reduced Hamiltonian after the mode-$1$ direction
has been eliminated: the amplitude $p$ (radial) and $s$ (tangential) of the slaved mode-$1$ displacement are
fixed by minimisation and by the constraint $J_1=0$. Mode $1$ is where Montaldi and Tokieda's $V_1'$ lies,
and their criterion uses the analogous reduced combination of $\delta$, $\gamma$ and $\beta$ (with Lemma~8's
weights, $\delta-3\gamma^2/\beta$; see Section~\ref{sec:remarks}). The difference is that $c_4$ is computed
on the level set $J^{-1}(\mu)$, so it also contains the curvature term $\tfrac14Q''(w_0)$ in the $J_1$
constraint, which does not enter a computation on the linear slice.
\end{remark}

\begin{remark}
The odd-$n$ derivation was done symbolically in $n$ and $\lambda$ by a second AI reader. It was then audited
step by step by a third reader: the coordinate change, the decoupling, completeness of the ansatz
(including the absence of aliasing at $n=5,7$), the slaved modes, the $J_1$ constraint, the residues and
Laurent degrees, and the invariant theory. The third reader also checked it by an independent exact route,
described next.
\end{remark}

\subsection{Exact values and independent verification}

Table~\ref{tab:c4} lists $c_4$ for $n=4,\dots,25$ from the closed forms. Every value has been checked
independently of the derivations above. The check is an order-by-order KKT expansion of the level-set
minimisation in \emph{Cartesian} coordinates, with all $2n$ coordinates unknown and no Fourier ansatz, no
symmetry reduction and no residue sums. It is carried out exactly over $\mathrm{GF}(p)$ for three primes
$p\approx2^{61}$ with $p\equiv1\pmod{4n}$, at both algebraic conjugates of $b_n$. In every case the
computation confirms three things exactly: the kernel eigenvalue vanishes, the cubic term vanishes, and
$c_4$ equals the tabulated rational number. A deliberately perturbed value is rejected (negative control).
Agreement modulo three primes of size $2^{61}$ is overwhelming evidence, though strictly it is not a proof
over $\mathbb Q$. The reduction from the printed coefficients to every table entry is, however, verified
exactly over $\mathbb Q$ in Lean (see Section~\ref{sec:limits}). The proofs are Theorems~\ref{thm:even} and~\ref{thm:odd}. The same engine in 40-digit
floating point, with the real root $b_n$, agrees to $10^{-37}$ or better. A separate 40-digit constrained
minimisation in the original chart gives the same values to 15 digits or better for
$n=4,\dots,13,15,17,19,21,25$; its linear-slice values are also shown.

\begin{table}[h]
\centering\small
\begin{tabular}{rrlrrl}
\toprule
$n$ & $b_n$ & geometry & $c_4$ (level set) & decimal & linear slice (numerical)\\
\midrule
4 & $0.2679491924$ & sphere & $1/4$ & $0.25$ & $0.2306624$\\
5 & $0.1715728753$ & sphere & $9/40$ & $0.225$ & $0.2046573$\\
6 & $0.0557280900$ & sphere & $23/30$ & $0.7666667$ & $0.8245485$\\
7 & $0$ & plane & $9/14$ & $0.6428571$ & $6/7=0.8571429$\\
8 & $-0.0627460668$ & $H^2$ & $95/56$ & $1.6964286$ & $2.2779413$\\
9 & $-0.1010205144$ & $H^2$ & $19/12$ & $1.5833333$ & $2.5781463$\\
10 & $-0.1428571429$ & $H^2$ & $49/15$ & $3.2666667$ & $4.9125$\\
11 & $-0.1715728753$ & $H^2$ & $75/22$ & $3.4090909$ & $5.7955980$\\
12 & $-0.2024108640$ & $H^2$ & $755/132$ & $5.7196970$ & $9.0145599$\\
13 & $-0.2251482266$ & $H^2$ & $675/104$ & $6.4903846$ & $10.9142245$\\
14 & $-0.2492775954$ & $H^2$ & $1693/182$ & $9.3021978$ & --\\
15 & $-0.2679491924$ & $H^2$ & $56/5$ & $11.2$ & $18.3284096$\\
16 & $-0.2875950927$ & $H^2$ & $1141/80$ & $14.2625$ & --\\
17 & $-0.3033370453$ & $H^2$ & $609/34$ & $17.9117647$ & $28.4269777$\\
18 & $-0.3197938296$ & $H^2$ & $3190/153$ & $20.8496732$ & --\\
19 & $-0.3333333333$ & $H^2$ & $27$ & $27$ & $41.5953947$\\
20 & $-0.3474164357$ & $H^2$ & $11139/380$ & $29.3131579$ & --\\
21 & $-0.3592455180$ & $H^2$ & $2175/56$ & $38.8392857$ & $58.2169269$\\
22 & $-0.3714998241$ & $H^2$ & $6145/154$ & $39.9025974$ & --\\
23 & $-0.3819660113$ & $H^2$ & $2475/46$ & $53.8043478$ & --\\
24 & $-0.3927725152$ & $H^2$ & $29183/552$ & $52.8677536$ & --\\
25 & $-0.4021298312$ & $H^2$ & $7227/100$ & $72.27$ & $103.3451777$\\
\bottomrule
\end{tabular}
\caption{Level-set effective quartic $c_4$ (Definition~\ref{def:c4}; $r_0=1$, $H$ as in~\eqref{eq:H},
Euclidean-unit kernel vector). All 22 rational values were confirmed by the exact $\mathrm{GF}(p)$
Cartesian KKT computation. The last column is the slaved quartic on the linear slice $x_0+N$, the
quantity that, on our reading, the computation of~\cite[\S3.5]{MT} evaluates (Section~\ref{sec:remarks}, P1
and P2). It
differs from $c_4$ in every case, and is positive in every case.}
\label{tab:c4}
\end{table}

%======================================================================
\section{Nonlinear stability for a degenerate minimum}\label{sec:nonlinear}
%======================================================================

The classical energy--momentum criteria assume a definite Hessian, a compact group, or both. Here the
Hessian is only semidefinite, and the group is $SE(2)$ or $PSL(2,\R)$ for $n\ge7$. We therefore give a
self-contained argument. It uses no stability theorem from the literature. The construction (a cross-section,
then an energy--momentum function with a penalty $(J_3-\mu_3)^2$) is classical in spirit, and we claim no
novelty for it.

\paragraph{Abstract setting.} Let $(P,\Omega)$ be a symplectic manifold and $H$ a smooth function. A Lie
group $G$ acts by symplectomorphisms preserving $H$. Functions $J_1,J_2,J_3$ are Hamiltonians for
generators $\xi_1,\xi_2,\xi_3$, where $\xi_3$ generates a circle action $K\cong SO(2)$, and
$\operatorname{span}(\xi_1,\xi_2)$ is $\operatorname{Ad}_K$-invariant. Let $x_0$ be a relative equilibrium,
$dF(x_0)=0$ for $F=H-\omega_0J_3$. Fix a $K$-invariant Riemannian metric on $P$. Assume:
\begin{itemize}
\item[(A1)] $J_1(x_0)=J_2(x_0)=0$, $K$ rotates the pair $(J_1,J_2)$, and $\{J_1,J_2\}(x_0)\ne0$.
\item[(A2)] $K$ acts freely near $x_0$.
\item[(A3)] There is a smooth curve $j\mapsto x_j$ of relative equilibria through $x_0$ with
$J_1=J_2=0$ and $J_3(x_j)=\mu_3+j$, with smooth angular velocities $\omega(j)$ and $\omega(0)=\omega_0$.
\item[(A4)] On $M_\mu=J^{-1}(\mu)/K$, the reduced Hamiltonian $h$ has positive semidefinite Hessian at
$[x_0]$, with kernel $K_0$ and positive definite on a complement $E$. The function $g(a)=h(a,\beta(a))-h(0)$
($\beta$ the slaving map) has vanishing cubic Taylor term and positive definite quartic Taylor term $q_4$.
\end{itemize}
Each $J_i$ is conserved, since $\{H,J_i\}=dH(\xi_i)=0$.

\begin{lemma}[Lemma A: cross-section]\label{lem:A}
Near $x_0$, $S:=\{J_1=J_2=0\}$ is a $K$-invariant symplectic submanifold of codimension $2$. It is invariant
under the flow of $H$, and the flow on $S$ is the Hamiltonian flow of $H|_S$. There is a $K$-invariant
neighbourhood $U$ of the whole orbit $K\cdot x_0$, and a continuous map $\gamma:U\to G$ with
$\gamma(x)\cdot x\in S$, $\gamma|_{K\cdot x_0}=e$ and $\gamma(kx)=k\gamma(x)k^{-1}$ for $k\in K$.
\end{lemma}

\begin{proof}
Since $\{J_1,J_2\}(x_0)\ne0$, $dJ_1$ and $dJ_2$ are independent. Moreover
$TS^\Omega=\operatorname{span}(X_{J_1},X_{J_2})$ with $\Omega(X_{J_1},X_{J_2})\ne0$, so $TS^\Omega$ and
$TS$ are complementary and $\Omega|_S$ is nondegenerate. Because $J_1,J_2$ are conserved, $X_H$ is tangent
to $S$, and $\Omega|_S(X_H,v)=dH(v)$ for $v\in TS$ gives $X_H|_S=X_{H|_S}$. $S$ is $K$-invariant because
$K$ rotates $(J_1,J_2)$.

For $\gamma$, let $\Psi(x,s_1,s_2)=(J_1,J_2)(\exp(s_1\xi_1+s_2\xi_2)\cdot x)$. At $(x_0,0)$, its derivative
in $(s_1,s_2)$ is the matrix $dJ_i(\xi_j(x_0))=\pm\{J_i,J_j\}(x_0)$, which is invertible by (A1). The
implicit function theorem gives $\varrho>0$, a neighbourhood $U_0$ of $x_0$, and a smooth $s(x)$ on $U_0$,
such that $s(x)$ is the \emph{unique} solution of $\Psi(x,s)=0$ with $|s|<\varrho$, and $s(x_0)=0$. Put
$\gamma(x)=\exp(s_1(x)\xi_1+s_2(x)\xi_2)$.

To extend to the whole orbit, note the following. For $k\in K$,
$k\exp(s\cdot\xi)k^{-1}=\exp(\operatorname{Ad}_k(s\cdot\xi))$ with $\operatorname{Ad}_k(s\cdot\xi)
\in\operatorname{span}(\xi_1,\xi_2)$, and $|\operatorname{Ad}_k s|=|s|$ for a $K$-invariant norm. Also
$S$ is $K$-invariant. So for $x\in kU_0$, the element $k\gamma(k^{-1}x)k^{-1}$ is the unique element
$\exp(s'\cdot\xi)$ with $|s'|<\varrho$ carrying $x$ into $S$. This uniqueness shows that the definitions on
overlapping translates $kU_0\cap k'U_0$ agree. Hence $\gamma$ is well defined and continuous on the
$K$-invariant neighbourhood $U=K\cdot U_0$, and it is $K$-equivariant.
\end{proof}

\begin{lemma}[Lemma B1: quartic lower bound on $M_\mu$]\label{lem:B1}
Under (A4), in coordinates $w=(a,b)\in K_0\times E$ centred at $[x_0]$, there are $\kappa,\varrho>0$ with
\[
h(a,b)-h(0)\ \ge\ \kappa\big(|a|^4+|b|^2\big)\qquad\text{for }|w|<\varrho .
\]
\end{lemma}

\begin{proof}
Let $f=h-h(0)$, so $d^2f(0)=\operatorname{diag}(0,B)$ with $B>0$. By the implicit function theorem there is a
smooth $\beta$ with $\partial_bf(a,\beta(a))=0$ and $\beta(0)=0$. Moreover $d\beta(0)=0$, because the cross
block vanishes, so $|\beta(a)|\le L|a|^2$. Taylor's theorem in $b$ about $\beta(a)$ gives
$f(a,b)-g(a)\ge\tfrac14\lambda_{\min}(B)|b-\beta(a)|^2$ for small $(a,b)$.

Next, $f(a,\beta(a))-f(a,0)=O(|a|^2)\cdot O(|a|^2)$, since $\partial_bf(a,0)=O(|a|^2)$ and $\beta=O(|a|^2)$.
So the Taylor terms of $g$ of order $\le3$ are those of $f|_{K_0}$. These vanish: order $2$ by (A4) and order
$3$ by assumption. Hence $g(a)\ge\tfrac12\min_{|a|=1}q_4\,|a|^4$ for small $a$. Finally
$|b|^2\le2|b-\beta|^2+2L^2|a|^4$.
\end{proof}

\begin{lemma}[Lemma B2: uniformity in nearby momentum]\label{lem:B2}
By (A2), $S/K$ is a manifold near $[x_0]$. Choose coordinates $(j,w)$ on it with $j=J_3-\mu_3$; this is
possible because $dJ_3\ne0$ on $S/K$, since $X_{J_3}\ne0$ and $\Omega|_S$ is nondegenerate. Choose them so
that $\{w=0\}$ is the curve $[x_j]$ of (A3), and $\{j=0\}=M_\mu$ with the coordinates of
Lemma~\ref{lem:B1}. Let $\Phi(j,w)=H(j,w)-H(j,0)$. Then there are $C,\varrho>0$ with
\[
\Phi(j,w)\ \ge\ \tfrac\kappa2\big(|a|^4+|b|^2\big)-Cj^2\qquad(|j|,|w|<\varrho).
\]
\end{lemma}

\begin{proof}
We have $\Phi(j,0)=0$. Also $\partial_w\Phi(j,0)=0$ for every $j$, because $x_j$ is a critical point of
$H$ on $\{J_3=\mu_3+j\}\cap S$. Differentiating these identities in $j$ gives $\partial_j\Phi(0,0)=0$ and
$\partial_w\partial_j\Phi(0,0)=0$. Hence $|\partial_j\Phi(0,w)|\le L'|w|^2=L'(|a|^2+|b|^2)$. Note that
there are no terms linear in $a$ or $b$ multiplied by $j$. Taylor's theorem in $j$ gives
\[
\Phi(j,w)\ge\Phi(0,w)-L'|j|\,(|a|^2+|b|^2)-C_2j^2 .
\]
For $\Phi(0,w)=h(w)-h(0)$ use Lemma~\ref{lem:B1}. The term $j|a|^2$ is handled by AM--GM: with
$X=\sqrt\kappa\,|a|^2$ and $Y=L'|j|/\sqrt\kappa$,
\[
L'|j|\,|a|^2=XY\le\tfrac12X^2+\tfrac12Y^2=\tfrac\kappa2|a|^4+\tfrac{L'^2}{2\kappa}j^2 .
\]
The term $j|b|^2$ is absorbed directly: $L'|j|\,|b|^2\le\tfrac\kappa2|b|^2$ once $|j|\le\kappa/(2L')$.
Altogether $\Phi\ge\tfrac\kappa2(|a|^4+|b|^2)-(L'^2/(2\kappa)+C_2)j^2$.
\end{proof}

\begin{remark}[The minimum does not persist at nearby momentum]\label{rem:nonpersist}
Lemma~\ref{lem:B2} only bounds how deep $\Phi$ can go below zero; it does not say $\Phi\ge0$. In fact the
strict minimum does \emph{not} persist. On the unstable side ($j<0$, slightly smaller ring), the kernel
eigenvalue of the reduced Hessian at $\mu_3+j$ is negative and of order $|j|$. The ring is then a saddle of
the reduced Hamiltonian, and the bifurcated minima sit at amplitude of order $|j|^{1/2}$ with depth of
order $j^2$.

We computed this numerically on the $D_{n/2}$-invariant family. For $n=8$, $\min\Phi/j^2$ equals
$-0.0022841$, $-0.0022845$ and $-0.0022845$ at $j=-10^{-2},-10^{-3},-10^{-4}$, and the minimiser satisfies
$\tau^*/\sqrt{|j|}\to0.1542$. For $n=12$, $\min\Phi/j^2\to-0.013679$ and $\tau^*/\sqrt{|j|}\to0.2007$. On
the stable side ($j>0$) the minimum is exactly $0$, attained at the ring.

So there is no neighbourhood on which the reduced Hamiltonians at all nearby momenta have uniform strict
minima. The step in Morse-lemma-based proofs that relies on such uniformity is therefore not available here. The
penalty function in Lemma~\ref{lem:B3} replaces it.
\end{remark}

\begin{lemma}[Lemma B3: conserved function with a strict minimum modulo $K$]\label{lem:B3}
On $S$ let $V=H-\omega_0J_3+K_c(J_3-\mu_3)^2$. Then $V$ is conserved and $K$-invariant. For $K_c$
large enough, the function it induces on $S/K$ has a strict local minimum at $[x_0]$.
\end{lemma}

\begin{proof}
Write $V(j,w)-V(0,0)=\Phi(j,w)+\psi(j)+K_cj^2$ with $\psi(j)=H(x_j)-H(x_0)-\omega_0j$. Along the family,
$d(H-\omega(j)J_3)(x_j)=0$ gives $\tfrac{d}{dj}H(x_j)=\omega(j)$. So $\psi(0)=\psi'(0)=0$ and
$|\psi(j)|\le C_3j^2$. By Lemma~\ref{lem:B2},
\[
V-V(0,0)\ge\tfrac\kappa2(|a|^4+|b|^2)+(K_c-C-C_3)j^2,
\]
which is positive for $(j,w)\ne0$ once $K_c>C+C_3$.
\end{proof}

\begin{lemma}[Lemma B4: Lyapunov stability on $S$]\label{lem:B4}
The orbit $K\cdot x_0$ is Lyapunov stable for the flow on $S$.
\end{lemma}

\begin{proof}
$K\cdot x_0$ is compact, and $V-V_0>0$ on a punctured $K$-invariant tubular neighbourhood $N_r$ of
$K\cdot x_0$ in $S$ (Lemma~\ref{lem:B3}). Given $\varepsilon<r$, let
$m=\min\{V-V_0:\ x\in S,\ \operatorname{dist}(x,K\cdot x_0)=\varepsilon\}$. This is a minimum over a
compact set, so $m>0$. Choose $\delta$ such that $V-V_0<m$ on the $\delta$-neighbourhood. Since $V$ is
conserved, a trajectory starting there never reaches distance $\varepsilon$. It stays in a compact,
collision-free set, so it exists for all time.
\end{proof}

\begin{theorem}\label{thm:abstract}
Under (A1)--(A4), $x_0$ is Lyapunov stable modulo $K$ in the full phase space $P$, with no restriction on
the momentum of the perturbation.
\end{theorem}

\begin{proof}
Let $x(0)$ be close to $K\cdot x_0$, put $\gamma=\gamma(x(0))$ (Lemma~\ref{lem:A}), and let
$y(t)=\gamma\cdot x(t)$. Since $G$ acts symplectically and preserves $H$, it commutes with the flow, so
$y(t)$ is the trajectory through $y(0)=\gamma\cdot x(0)\in S$. Moreover $y(0)\to K\cdot x_0$ and
$\gamma\to e$ as $x(0)\to K\cdot x_0$. By Lemma~\ref{lem:B4}, $y(t)$ stays within $\varepsilon'$ of
$K\cdot x_0$. Then
\[
\operatorname{dist}(x(t),K\cdot x_0)\le d(\gamma^{-1}y(t),y(t))+\operatorname{dist}(y(t),K\cdot x_0)
\le\omega(\gamma)+\varepsilon',
\]
where $\omega(\gamma)=\sup\{d(\gamma^{-1}y,y)\}$ over a fixed compact neighbourhood of $K\cdot x_0$. This
tends to $0$ as $\gamma\to e$ by continuity of the action.
\end{proof}

\paragraph{Proof of Theorem~\ref{thm:main}.} We verify (A1)--(A4) for the ring at $b_n$, $n\ge4$, in all
three geometries.
\begin{itemize}
\item (A1): $J_+(x_0)=0$ by $\Z_n$ symmetry, and a rotation by $\psi$ multiplies $J_+$ by $e^{i\psi}$.
$\{J_1,J_2\}(x_0)=n(1-b_n)/(2(1+b_n))\ne0$ by~\eqref{eq:bracket}, since $|b_n|<1$. The span of
$\xi_1,\xi_2$ is $\operatorname{Ad}_K$-invariant in $\mathfrak{so}(3)$, $\mathfrak{se}(2)$ and
$\mathfrak{sl}(2,\R)$.
\item (A2): a nontrivial rotation moves a labelled ring.
\item (A3): the concentric rings of radius $r$ are relative equilibria for every $r$~\cite[(2.1)]{MT}, with
$J_+=0$, and $J_3=nr^2/(1+\lambda r^2)$ is strictly increasing in $r$.
\item (A4): Proposition~\ref{prop:spectral} gives the Hessian and its kernel. By the symmetry argument after
Definition~\ref{def:c4}, $q_4=c_4|a|^4$ in a chart, and $c_4>0$ by Theorems~\ref{thm:even}
and~\ref{thm:odd}.
\end{itemize}
Theorem~\ref{thm:abstract} now gives Theorem~\ref{thm:main}. \qed

%======================================================================
\section{Relation to the literature}\label{sec:lit}
%======================================================================

\paragraph{Patrick--Roberts--Wulff.} The following is an independent confirmation for the sphere and the
hyperbolic plane. Patrick, Roberts and Wulff~\cite{PRW} consider a free action with
$\operatorname{Ad}^*$-equivariant momentum map. Their Theorem~4.2 states that $p_e$ is $G$-stable if
$H^{-1}(H(p_e))\cap J^{-1}(T_2^{U_{\mu_e}}(\mu_e))\cap U_S=\{p_e\}$ for some neighbourhoods $U_S$ (in a slice)
and $U_{\mu_e}$. No Hessian condition is required. When $\mathfrak g^*/G$ is Hausdorff at $\mu_e$, which
holds if $\mu_e$ is regular~\cite[Prop.~4.10(1)]{PRW}, this condition is equivalent to $p_e$ being isolated
in its energy level within the reduced space. A strict local minimum (Lemma~\ref{lem:B1}) is isolated.
Their Corollary~5.4 then upgrades $G$-stability to $G^o_{\mu_e}$-stability whenever there is a
$G^o_{\mu_e}$-invariant inner product on $\mathfrak g^*$, which is automatic here since $G_\mu=SO(2)$ is
compact.

For $SO(3)$ ($\mu\ne0$) and $PSL(2,\R)$ ($\mu$ timelike), $\mu$ is regular, and the action is free and
proper on a neighbourhood of the labelled ring. With $J$ shifted to its equivariant form (for
$\lambda\ne0$), these results therefore give $G_\mu$-stability of the ring at $b_n$ for $n\ne7$, once the
strict minimum of the reduced Hamiltonian is known, i.e.\ once Sections~\ref{sec:setup}
and~\ref{sec:quartic} are in hand.

\paragraph{The plane.} For $\lambda=0$ the momentum map of $SE(2)$ is equivariant only for an affine
action (equation~\eqref{eq:bracket}: $\{J_1,J_2\}=n/2$ is a cocycle). In the central extension that makes
it equivariant, the centre acts trivially, so the action is not free. So~\cite{PRW} does not apply
verbatim, and the self-contained Theorem~\ref{thm:abstract} is what covers $n=7$.

\paragraph{Montaldi 1997.} Montaldi--Tokieda's Proposition~9 relies on~\cite{Montaldi97} (cited as
``cf.~[17]'' in its proof~\cite[p.~20]{MT}, and as ``[17, Theorem~1.2]'' on p.~10), whose setting
is a compact group. For the sphere it applies; in the degenerate case we use it on the reduced
Hamiltonian, i.e.\ with the level-set quartic. We have read only that paper's abstract and the way
it is quoted in~\cite[pp.~10, 20]{MT}; we rely on it nowhere. Montaldi and Rodr\'\i guez-Olmos~\cite{MRO} treat
noncompact groups, but under formal stability (a definite Hessian on the symplectic normal space), so the
degenerate case is not covered verbatim. We did not find a published statement that covers $SE(2)$,
$PSL(2,\R)$ and the degenerate case together, which is why Section~\ref{sec:nonlinear} gives a self-contained
argument.

%======================================================================
\section{Comparison with intermediate steps of the published argument}\label{sec:remarks}

Montaldi and Tokieda's conclusions (stability at $b_n$ for even $n$ and for $n=3,5,7,9,11$) are correct,
and our results confirm them. Their framework (the deformation family, the Fourier blocks, the identification
of $b_n$ and the critical modes, the invariant-theoretic Proposition~12) is what makes the present proof
possible. The notes below compare our computation with intermediate steps of~\cite[\S3.5]{MT} only. They
reflect our reading of the published text, and we may have missed conventions the authors intended. We give
page and equation numbers of~\cite{MT} (arXiv:1210.5662v1), so that each point can be checked line by line. The scripts are
in \code{conj13/review/}.

\paragraph{Normalisation dictionary.} Montaldi and Tokieda's augmented Hamiltonian is
$\widehat H_\lambda=H_\lambda-\omega\sum_j\kappa|z_j|^2/\sigma_j$ with
$H_\lambda=-\frac{\kappa^2}{4\pi}\sum_{i<j}\log\frac{|z_i-z_j|^2}{\sigma_i\sigma_j}$~\cite[p.~8]{MT}.
With $\kappa=1$, $r_0=1$ and $\omega=\omega_0$ of~\cite[(2.1)]{MT}, this is $\widehat H=F/(4\pi)$, where
$F=H-\omega_0J_3$ is as in Section~\ref{sec:setup}. Their Fourier vectors~\cite[(2.5)]{MT} are
$\alpha_r^{(\ell)}=\sum_j\cos(2\pi\ell j/n)\,\delta r_j$ (and similarly $\beta_r^{(\ell)}$,
$\alpha_\theta^{(\ell)}$, $\beta_\theta^{(\ell)}$), with $\delta r_j$ the unit radial vector at vortex $j$, so
they have Euclidean length $(n/2)^{1/2}$ for $0<\ell<n/2$. On $V_c=\operatorname{Ker}B$~\cite[p.~22]{MT},
$z=x+iy$ are coordinates with respect to $\alpha_r^{(l)},\beta_r^{(l)}$, $l=(n-1)/2$; hence
$x=(2/n)^{1/2}\varepsilon$, with $\varepsilon$ our coordinate along $e$ (Definition~\ref{def:c4}). On $V_1'$,
$w=u+iv$ are coordinates with respect to the basis~\cite[(2.8)]{MT}, whose first vector is
$v_1=\sigma\alpha_r^{(1)}-\widetilde\sigma\beta_\theta^{(1)}$, with $\sigma=1+\lambda r_0^2$ and
$\widetilde\sigma=1-\lambda r_0^2$. With $B=d^2\widehat H$, $C=d^3\widehat H$, $D=d^4\widehat H$ at the ring, as
in Lemma~8~\cite[p.~19]{MT}, and $\alpha=\alpha_r^{(l)}$,
\begin{equation}\label{eq:dict}
\begin{gathered}
\beta=B(v_1,v_1)=\frac{n(n-1)\widetilde\sigma^2}{4\pi},\qquad
\gamma=C(\alpha,\alpha,v_1)=\frac{\partial^3\widehat H}{\partial x^2\partial u},\\
\delta=D(\alpha,\alpha,\alpha,\alpha)=\frac{24}{4\pi}\Big(\frac n2\Big)^2q_4=\frac{3n^2}{2\pi}\,q_4 .
\end{gathered}
\end{equation}
The formula for $\beta$ is~\cite[(2.10) and p.~23]{MT}. Here $q_4$ is the $\varepsilon^4$ coefficient of $F$
along the straight line $x_0+\varepsilon e$, with no slaving (the \emph{unslaved quartic}; $q_4=51/28$ at
$n=7$). We checked~\eqref{eq:dict} by evaluating $B$, $C$, $D$ directly in Montaldi and Tokieda's
normalisation (60-digit finite differences, \code{review/p2_mt_coeffs.py}). At $n=7$ this gives
$\beta=21/(2\pi)$, $\gamma=63/(4\pi)$ and $\delta=1071/(8\pi)$, exactly the printed values~\cite[p.~23]{MT}. At
$n=9$ it gives $6.946$, $13.839$, $182.435$ (printed: $6.9$, $13.8$, $182.4$), and at $n=11$ it gives
$12.015$, $29.704$, $555.141$ (printed: $12.0$, $29.7$, $555.1$). So the printed $\gamma$ is $C(\alpha,\alpha,v_1)$,
and the trilinear term is $C(a,a,b)=\gamma\operatorname{Re}(z^2w)$ for $a=x\alpha+y\beta_r^{(l)}$ and
$b=uv_1+vv_2$. (The text before~(3.9) writes $Ca^2b=\gamma(z^2w+\bar z^2\bar w)$ and
$\gamma=\tfrac14\,\partial^3(Ca^2b)/\partial x^2\partial u=\tfrac32\,\partial^3\widehat H/\partial x^2\partial u$.
These normalise $\gamma$ differently from the printed numbers, by constant factors. The ratio test below uses
only the printed numbers.)

With~\eqref{eq:dict}, Lemma~8's quantity for $a\in V_c$ and $b\in V_1'$ is
\begin{equation}\label{eq:lemma8}
Da^4+6Ca^2b+3Bb^2=\delta|z|^4+6\gamma\operatorname{Re}(z^2w)+3\beta|w|^2,\qquad
\min_w(\cdots)=\Big(\delta-\frac{3\gamma^2}{\beta}\Big)|z|^4 .
\end{equation}
This is $24$ times the minimum over $w$ of the Taylor polynomial
$\tfrac1{24}\delta|z|^4+\tfrac12\gamma\operatorname{Re}(z^2w)+\tfrac12\beta|w|^2$ of $\widehat H$ on $N$. In our
units, therefore,
\begin{equation}\label{eq:dictlin}
\delta-\frac{3\gamma^2}{\beta}=\frac{3n^2}{2\pi}\,c_4^{\rm lin},
\end{equation}
where $c_4^{\rm lin}$ is the slaved quartic on the linear slice $x_0+N$ (last column of Table~\ref{tab:c4}).

\paragraph{P1: linear slice versus level set.} Proposition~9~\cite[p.~20]{MT} applies Lemma~8 ``with
$\mathbb R^n$ replaced by $N$ and $f$ by $\widehat H$''. Its proof ends: ``The claim now follows because the
hypotheses of Lemma~8 applied on $M_\mu$ are equivalent to the same hypotheses applied on $N$.'' The
equivalence is clear for the second-order hypothesis ($B\ge0$, with the same kernel). For the quartic
hypothesis with degenerate $B$ it seems to us to need a further argument, for the reason given in
Section~\ref{sec:quartic}. A curve in
$M_\mu$ lifts to a curve in $J^{-1}(\mu)$ tangent to $N$, but its second-order part leaves $N$, chiefly along
the radial part of $V_0$, and that component couples to $a^2$ at cubic order. Lemma~10~\cite[p.~20]{MT} and
Proposition~12(iii)~\cite[p.~22]{MT} set the coupling to $V_0$ aside because $V_0\cap N=\{0\}$. That is
right on $N$; on our reading, it also means that the computation evaluates the quartic on the linear slice. For even $n$, Lemma~10 makes the
straight-line $t^4$ coefficient of~\cite[p.~21]{MT} the linear-slice quartic. For odd $n$, the printed
numbers evaluate the linear-slice quantity, by~\eqref{eq:dictlin} and the ratio test below. The two
quantities differ:
\[
n=7:\ \ c_4=\tfrac9{14}\ \text{(level set)}\ \text{vs.}\ c_4^{\rm lin}=\tfrac67;\qquad
n=8:\ \ c_4=\tfrac{95}{56}=1.6964\ \text{vs.}\ c_4^{\rm lin}=2.2779 .
\]
Table~\ref{tab:c4} shows the difference for every $n$ computed. The linear-slice value is smaller for
$n=4,5$ and larger for $n\ge6$, so in general the two need not have the same sign. For the vortex ring both quantities are positive at every $n$ in
Table~\ref{tab:c4}, and $c_4>0$ for all $n$, so no conclusion of~\cite{MT} changes.

\paragraph{P2: the weights of Lemma~8 in~(3.9).} Lemma~8~\cite[p.~19]{MT} requires $Ca^3=0$ and
$Da^4+6Ca^2b+3Bb^2>0$. On p.~23, ``the key quantity $Da^4+Ca^2b+Bb^2$ becomes, on completing the square,''
\[
\delta|z|^4+2\gamma\operatorname{Re}(z^2w)+\beta|w|^2
=\delta\Big|z^2+\frac\gamma\delta\bar w\Big|^2+\frac{\beta\delta-\gamma^2}{\delta}|w|^2
\qquad\text{\cite[(3.9)]{MT}},
\]
and the criterion is stated as $\beta\delta>\gamma^2$. As displayed, this combination does not carry the
weights $6$ and $3$ of Lemma~8; we take them to be intended. With the dictionary above, the weighted quantity
is~\eqref{eq:lemma8}, and the corresponding criterion is $\delta-3\gamma^2/\beta>0$, i.e.\
$\beta\delta>3\gamma^2$.

\emph{Ratio test.} The ratio $(\delta-k\gamma^2/\beta)/\delta$ does not change if $\widehat H$ or the bases of
$V_c$ and $V_1'$ are rescaled. It can therefore be compared directly with our $c_4^{\rm lin}/q_4$ and
$c_4/q_4$, which are computed independently of~\cite{MT}:
\begin{center}\small
\begin{tabular}{rcccc}
\toprule
$n$ & $(\delta-\gamma^2/\beta)/\delta$ & $(\delta-3\gamma^2/\beta)/\delta$ & $c_4^{\rm lin}/q_4$ &
$c_4/q_4$ (level set)\\
 & \multicolumn{2}{c}{(printed $\beta,\gamma,\delta$ of~\cite[p.~23]{MT})} & (ours) & (ours)\\
\midrule
7 & $14/17\approx0.8235$ & $8/17\approx0.4706$ & $8/17\approx0.47059$ & $6/17\approx0.35294$\\
9 & $0.8487$ & $0.5461$ & $0.54654$ & $0.33565$\\
11 & $0.8676$ & $0.6027$ & $0.60315$ & $0.35478$\\
\bottomrule
\end{tabular}
\end{center}
At $n=9,11$ the printed values have three or four significant figures. Script:
\code{review/p2_ratio_test.py}, output \code{review/p2_ratio_test.out}. With Lemma~8's weights, the printed
numbers reproduce our linear-slice ratio exactly at $n=7$, and to the printed precision at $n=9$ and $11$.
The unweighted combination matches neither of our columns, which is why we read the weights as intended.
Three points of detail follow.

(a) Montaldi and Tokieda do not print the number $6/7$. It is our computation of $c_4^{\rm lin}$ at $n=7$.
Their printed values reproduce it through~\eqref{eq:dictlin}:
$\delta-3\gamma^2/\beta=\frac{1071}{8\pi}-3\big(\frac{63}{4\pi}\big)^2\big/\frac{21}{2\pi}=\frac{63}\pi
=\frac{3\cdot49}{2\pi}\cdot\frac67$.

(b) An alternative reading. At $n=7$, $\delta-\gamma^2/\beta=441/(4\pi)$, and $441/4=\frac{7^3}2\cdot\frac9{14}$.
This suggests that the unweighted combination might be the level-set value $c_4=9/14$ under a scale
$n^3/(2\pi)$. The numbers do not support this reading. The scale is fixed by $\delta$: by~\eqref{eq:dict}, $\delta/q_4=3n^2/(2\pi)=147/(2\pi)$
at $n=7$, not $343/(2\pi)$. Under that scale $\delta-\gamma^2/\beta$ corresponds to $3/2$, not $9/14$. In
normalisation-free form, $(\delta-\gamma^2/\beta)/\delta=14/17$, while $c_4/q_4=6/17$. Nor does the reading
fit at $n=9$: there $\delta-\gamma^2/\beta=182.4-13.8^2/6.9=154.8$, whereas $\frac{n^3}{2\pi}c_4=\frac{729}{2\pi}
\cdot\frac{19}{12}\approx183.7$ (and $\frac{3n^2}{2\pi}c_4\approx61.2$). The agreement at $n=7$ appears to be a numerical coincidence.

(c) The conclusions are unaffected. With the weights, the printed numbers give $\beta\delta-3\gamma^2=661.5/\pi^2$
($n=7$), about $687$ ($n=9$) and about $4015$ ($n=11$), in the units of~\cite{MT}. All are positive,
consistent with $c_4^{\rm lin}>0$ in Table~\ref{tab:c4}. In any case Theorem~\ref{thm:odd} gives $c_4>0$ on the
level set.

\paragraph{The printed values at $n=5$.} At $n=5$~\cite[p.~23]{MT} print $\beta=37.1/r_0^2$,
$\gamma=-9.04/r_0^3$ and $\delta=15.8/r_0^4$. The formula $\beta=n(n-1)\kappa^2\widetilde\sigma^2/(4\pi r_0^2)$ on
the same page gives $\beta\approx1.092/r_0^2$ at $n=5$ ($\kappa=1$, $\lambda r_0^2=b_5=3-2\sqrt2$), whereas at
$n=7,9,11$ it reproduces the printed $\beta$. Our direct evaluation at $n=5$ gives $\beta\approx1.092$,
$|\gamma|\approx1.085$ and $\delta\approx5.679$; with these, $(\delta-3\gamma^2/\beta)/\delta=0.4302=c_4^{\rm lin}/q_4$.
The printed values give $0.5818$ (weighted) and $0.8606$ (unweighted). These match neither
$c_4^{\rm lin}/q_4=0.4302$ nor $c_4/q_4=0.4730$. We have not been able to reconcile the printed $n=5$
values with these formulas, and we may be missing a convention specific to that case; we do not use them. The conclusion for $n=5$ does not depend on them: Theorem~\ref{thm:odd} gives
$c_4=9/40>0$, and $c_4^{\rm lin}\approx0.2047>0$ as well.

\paragraph{P3: the printed even-$n$ bracket $T$.} We recomputed symbolically the straight-line
expansion $z_j=(r_0+(-1)^jt)e^{2\pi ij/n}$ of~\cite[\S3.5.1, p.~21]{MT}. Our $t^4$ coefficient agrees with an
independent 40-digit evaluation to 15 digits, but differs from the printed bracket $T$ (p.~21) by a small
relative amount: $+0.10\%$ at $n=8$, $+0.49\%$ at $n=10$, $+0.06\%$ at $n=6$ and $-1.8\%$ at $n=4$ (printed
value relative to ours). The differences are small and may come from rounding or transcription. They change no sign, and the even-$n$
conclusion is in any case re-proved by Theorem~\ref{thm:even}.

%======================================================================
\section{Numerical confirmation}\label{sec:numerics}
%======================================================================

We integrated the $n=8$ ring in the hyperbolic plane directly. This is a check on the theory, not part of
the proof. The code is independent of that used for $c_4$. It uses the Poincar\'e disk with area form
$4\,dx\wedge dy/(1-|z|^2)^2$, DOP853 with relative tolerance $10^{-12}$, and final time $T=2000$, which is
more than six periods of the nonlinear oscillation at amplitude $10^{-4}$.

The initial data are the ring at $\xi=\xi^*(1+\mathrm{rel})$, where $\xi=|z|^2$ and $\xi^*=-b_8$ is the
threshold, plus the critical mode $|z_k|=\sqrt\xi(1+\delta(-1)^k)$. Here $\mathrm{rel}>0$ is the stable
side. The diagnostic is $D(t)=\max_{i<j}|d_{ij}(t)-d_{ij}^{\rm ring}|$ in hyperbolic distance, which is
invariant under the whole group. Energy and momentum are conserved to $10^{-8}$ and $6\cdot10^{-10}$, and the
unperturbed ring drifts by less than $6\cdot10^{-12}$.

\begin{center}\small
\begin{tabular}{lrrrrr}
\toprule
$\max_tD(t)/D(0)$ & rel $=-10^{-2}$ & $-10^{-3}$ & $0$ & $+10^{-3}$ & $+10^{-2}$\\
\midrule
$\delta=10^{-3}$ & 29.34 & 9.13 & 1.000 & 1.000 & 1.000\\
$\delta=10^{-4}$ & 293.6 & 91.86 & 1.000 & 1.000 & 1.000\\
\bottomrule
\end{tabular}
\end{center}

Below threshold the deviation saturates at an amplitude independent of $\delta$: about $0.029$ for
rel $=-10^{-2}$ and $0.0091$ for rel $=-10^{-3}$. Their ratio is close to $\sqrt{10}$. This is the
$\sqrt{\Delta}$ signature of a supercritical pitchfork: the orbit wanders near the stable bifurcated pair of
$n/2$-gons. At threshold and above it, $D(t)$ never exceeds $D(0)$. It oscillates, spends $38$--$42\%$ of
the time below $D(0)/2$ at threshold, and shows no secular growth. If the quartic were negative, threshold
runs would leave the $\delta$-neighbourhood on the nonlinear time scale. The same behaviour was observed for
the image (``tanh'') hyperbolic model, which is not part of Conjecture~13.

%======================================================================
\section{Limitations, status and computer-assisted parts}\label{sec:limits}
%======================================================================

\paragraph{Review status.} This draft has been checked only by independent AI readers (separate model
instances that re-derived the results with their own code before reading the candidate argument) and by
computer calculations. \textbf{No human specialist has reviewed it.} Its claims should be treated
accordingly until such a review has taken place.

\paragraph{Literature not read.} We have not read Kurakin--Yudovich~\cite{KY} beyond its abstract, nor
Kurakin~\cite{Kurakin03,Kurakin04}. We make no claim that our results agree with, extend, or differ from
theirs. One unverified observation: if $H=-\sum\log|z_i-z_j|$ and a per-vortex radial amplitude are used,
our level-set value at $n=7$ rescales to $\tfrac9{14}\cdot\tfrac{49}8=\tfrac{63}{16}$ (the linear-slice
value would rescale to $21/4$). The number $63/16$ has been attributed to~\cite{KY} as the heptagon's
quartic coefficient in secondary notes we started from, but we have not verified that attribution or the
normalisation used there. This is therefore an unverified consistency remark, not a claim of agreement. Kurakin's sphere results may
already cover $n=4,5,6$ at the critical latitude.

\paragraph{What is computer-assisted.} All computations are exact symbolic (sympy) or exact modular unless
labelled numerical. Paths are relative to \code{conj13/} in the accompanying repository.
\begin{itemize}
\item \emph{Even-$n$ closed form} (Theorem~\ref{thm:even}). The expansion of $H$ on
$\operatorname{Fix}(D_M)\cap J^{-1}(\mu)$, symbolic in $M$ and $\lambda$, and its reduction modulo~\eqref{eq:bn}
were done twice, independently: \code{even_exact.py}, \code{even_formula.py} and
\code{review/even_check.py}. The product formulas were checked by resultants for $M\le6$.
\item \emph{Odd-$n$ closed form} (Theorem~\ref{thm:odd}). The symbolic derivation in $n$ and $\lambda$,
including the residue evaluation of root-of-unity sums with Laurent-degree checks, is in
\code{review/odd_general.py}, \code{review/odd_general_run.py} and \code{review/odd_assemble.py}. The
human-readable steps are reproduced above. The output coefficients $f_2$, $E_{ij}$ and $a_2$ are printed in
the proof so that they can be re-derived by hand or with another system.
\item \emph{Independent exact check of Table~\ref{tab:c4}.} Cartesian order-by-order KKT over
$\mathrm{GF}(p)$, with no ansatz: \code{review3/kkt_exact.py} and \code{review3/run_exact.py} (odd $n$ and
$n=8$), and \code{paper/table_check.py} (all $n=4,\dots,25$, with a negative control; output in
\code{paper/table_check.out}).
\item \emph{Spectral checks} (Proposition~\ref{prop:spectral}, numerical, 40 digits):
\code{gaps/spectral_check.py}, plus the bordered inertia in \code{review3/hess_check.py}.
\item \emph{Lemma~\ref{lem:B3} sanity check} (numerical): \code{gaps/review/indep_check.py}. For
$n=4,\dots,9$, the minimum of $V-V_0$ on spheres of radius $\varepsilon$ in $S$ scales as $c_4\varepsilon^4$.
\item \emph{Remark~\ref{rem:nonpersist}}: \code{gaps/nearby_momentum.py}.
\item \emph{P1--P3}: \code{engine.py}, \code{reduced_quartic*.py}, \code{review/c4_levelset.py},
\code{review/p2_check.py}, \code{review/p2_ratio_test.py} (the ratio table in P2),
\code{review/p2_mt_coeffs.py} (direct evaluation of $\beta,\gamma,\delta$ in the normalisation of~\cite{MT}),
\code{review/p3_check.py}, \code{mt_n7_check.py}.
\item \emph{Integrations}: \code{integrate.py} with outputs in \code{out/}.
\end{itemize}
The parts of the proof that use a computer are the algebraic simplifications in Theorems~\ref{thm:even}
and~\ref{thm:odd}: series expansion, residue sums, and reduction modulo a quadratic relation. Each was
reproduced by an independent route. The structural steps (Proposition~\ref{prop:spectral},
Section~\ref{sec:nonlinear}) are pen-and-paper.

\paragraph{Formal verification (Lean 4).} The exact-algebra layer is formalised in Lean~4 with Mathlib,
in \code{lean/Havelock/Conj13Algebra.lean}, with no unproved steps and only the standard axioms. The
repository's continuous-integration build checks the Lean files. The precise scope is as follows. The
coefficients printed in the proofs of Theorems~\ref{thm:even} and~\ref{thm:odd}, namely $h_2,h_4$ (even $n$)
and $f_2$, $E_{00},E_{10},E_{01},E_{20},E_{02}$, $a_2$ (odd $n$), are explicit closed-form rational
expressions in $n$ (or $M=n/2$) and $b=\lambda r_0^2$. Lean takes these expressions as definitions and proves,
symbolically in $n$ and $b$, everything downstream of them:
\begin{itemize}
\item that $h_2$ and $f_2$ vanish exactly on the threshold relation~\eqref{eq:bn};
\item that eliminating $s$ by the $J_1$ constraint and minimising over $p$ (odd $n$), and reducing modulo the
threshold relation (both parities), gives the stated closed forms for $c_4$, for every $n$;
\item that these equal $c_4$ at the true (irrational) $b_n$ and are positive for every $n\ge4$, and every entry
of Table~\ref{tab:c4}, exactly over $\mathbb Q$.
\end{itemize}
It also proves the threshold facts ($b_n$ in closed form, its uniqueness in $(-1,1]$ and its sign pattern,
for all $n\ge4$), the spectral algebra of Proposition~\ref{prop:spectral} (including the Fej\'er identity
for all $n$ and the eigenvalues $3/2$ and $8/3$ at $n=4,5$), the comparisons in P1, and the
completing-the-square identities behind P2. What Lean does \emph{not} check is the derivation of the closed
forms $h_2,h_4,f_2,E_{ij},a_2$ from the Hamiltonian: the series expansions along the symmetric curve (even
$n$) and along the ansatz (odd $n$), and the root-of-unity residue sums for general $n$. Nor does it check the
normalisation factors relating $\tau$ and $A$ to $\varepsilon$. These steps are computer algebra, reproduced
independently and checked against the exact $\mathrm{GF}(p)$ Cartesian computation as described above.
Section~\ref{sec:nonlinear} is not formalised.

\paragraph{Scope.} The result concerns the family~\eqref{eq:H} of~\cite{MT}, i.e.\ the ``sinh'' model in
the hyperbolic case. It says nothing about the image (``tanh'') hyperbolic model or the flat disk with a
wall. The stability obtained is Lyapunov stability modulo $SO(2)$. We do not address stability on
exponentially long times for perturbations of the full system, KAM-type questions, or the bifurcated
branches beyond what~\cite[Remark~11]{MT} states.

\section*{Acknowledgements and disclosure of AI assistance}
This work was developed with substantial assistance from an AI system (Claude, Anthropic). The AI
assistance covered the derivations, the computer-algebra and numerical code, the independent re-derivations
used as ``readers'', and the drafting of this manuscript. The author directed the work and takes
responsibility for the content. As stated in Section~\ref{sec:limits}, no human expert has yet checked the
mathematics, and we welcome corrections. We thank James Montaldi and Tadashi Tokieda, whose paper posed
the problem and supplied the framework used here.

\bibliographystyle{plain}
\bibliography{refs}

\end{document}
